% Generated draft: requires mathematical and source review.
\documentclass[11pt]{article}
\usepackage[margin=25mm]{geometry}
\usepackage{fontspec}
\setmainfont{Latin Modern Roman}
\usepackage{amsmath,amssymb,mathtools,unicode-math}
\setmathfont{Latin Modern Math}
\usepackage{enumitem,needspace,xurl,hyperref}
\setlength{\parindent}{0pt}
\setlength{\parskip}{5pt}
\setlength{\emergencystretch}{3em}
\newcommand{\N}{\mathbb{N}}
\newcommand{\Z}{\mathbb{Z}}
\newcommand{\Q}{\mathbb{Q}}
\newcommand{\R}{\mathbb{R}}
\newcommand{\C}{\mathbb{C}}
\newcommand{\F}{\mathbb{F}}
\newcommand{\A}{\mathbb{A}}
\newcommand{\PP}{\mathbb{P}}
\newcommand{\E}{\mathbb{E}}
\newcommand{\eps}{\varepsilon}
\newcommand{\dd}{\,\mathrm{d}}
\begin{document}
\section*{UnsolvedMath ID: 20002874}
\textbf{Source problem code:} catalogue entry
\subsection{Definitions and mathematical statement}
Let \(M\subset\C^n\), \(n\ge2\), be a compact smooth CR manifold given by real polynomial equations in \(z,\bar z\). Its CR dimension is the complex dimension of its complex tangent bundle, and its CR codimension is \(\dim_{\R}M-2\dim_{\rm CR}M\), assumed at least one. For a chosen polynomial presentation \(M=\{P_1=\cdots=P_s=0\}\), record its maximum degree \(D=\max_j\deg P_j\). Suppose, under appropriate filling hypotheses, that an associated Harvey--Lawson holomorphic variety or chain \(V\) with boundary \(M\) exists and has isolated singularities. The source question is to relate the number
\[
\nu(V)=\#\operatorname{Sing}V
\]
to the degrees of the defining polynomials of \(M\), specifying any necessary dependence on dimension and presentation complexity.

The existence and uniqueness hypotheses for such a filling, the meaning of the associated variety when chains or reducible fillings occur, and whether singularities are counted with multiplicities are not supplied. In particular, the CR codimension assumption alone does not specify the maximal-complexity and other hypotheses needed for a holomorphic boundary theorem. A polynomial degree is presentation-dependent: replacing \(P_j\) by \(hP_j\) with a strictly positive real polynomial \(h\) preserves its real zero set while increasing degree. This does not exclude upper bounds for fixed presentations, but prevents treating presentation degree as an intrinsic invariant without normalization.

The supplied research aid gives a conditional bound in terms of filling degrees instead. Let \(X\subset\PP^{p+c}_{\C}\), \(p\ge1\), be a reduced projective complete intersection of dimension \(p\) and multidegree \((e_1,\ldots,e_c)\), with isolated projective singular scheme \(\Sigma\). Here \(\Sigma\) is cut out on \(X\) by the maximal minors of the Jacobian matrix of the defining equations. The reported Jacobian--B\'ezout estimate is
\[
\#\operatorname{Sing}X\le\operatorname{length}(\Sigma)
\le\left(\prod_{i=1}^ce_i\right)\left(\sum_{i=1}^c(e_i-1)\right)^p.
\]
For a filling contained in such an \(X\), this bounds its interior singular points when they are singular points of \(X\). It requires isolated singularities of the entire projective singular scheme, including at infinity; isolated singularities in an affine piece alone do not supply that hypothesis.

The aid also reports, for each \(d\ge2\), a degree-\(d\) polynomial CR boundary in \(\C^2\) with a unique reducible curve filling having exactly \(\binom d2\) ordinary nodes. The construction and the precise class in which uniqueness holds are not given here. This family indicates that quadratic growth is already required in that curve-filling setting. The remaining transfer problem is an effective estimate of filling multidegrees from a specified, normalized measure of boundary complexity. The conditional projective estimate alone is not such an estimate and does not settle the original general question.
\subsection{Short English statement}
Relate algebraic complexity of a compact CR boundary to the isolated singularities of its holomorphic filling. The supplied partial result bounds singularities using projective complete-intersection filling degrees and reports a quadratic lower family, while the boundary-to-filling complexity step and filling hypotheses remain unspecified.
\subsection{Sources}
\begin{itemize}
\item S. Ji, Problem 3 in the AIM workshop \emph{Emerging applications of complexity for CR mappings}, \url{https://aimath.org/WWN/crmappings/crmappings.pdf}. The source names the Harvey--Lawson variety but gives no full bibliography for the filling theory.
\item Local UnsolvedMath registry, permanent ID 20002874, exact code \texttt{catalogue entry}; canonical location \texttt{aim-several-complex-variables-notes.json}, notes[31], dated 14 August 2026. The statement is labeled exact; its verification reports no substantive OCR error.
\item Research attempt 1: partial result, problem open at the run, rigor label proved, novelty confidence low; registry status partially solved. The projective bound and node family are supplied by a machine-generated research aid, not peer reviewed, and require independent verification before citation.
\end{itemize}
\end{document}
